% Capítulo 1 — Introdução. Reescrita em 2026-11-16 pela fórmula de Head (via
% Bellemare, Doing Economics, cap. 2; docs/framework-dissertacao.md §3.1):
% gancho no mundo real, pergunta como primeira frase, antecedentes como
% história, contribuição, organização. Promete só o que o corpo entrega.

\chapter{Introdução}
\label{cap:introducao}

Em abril de 2010, José Maria Filho, o Zé Maria do Tomé, foi assassinado na
Chapada do Apodi, no interior do Ceará. Ele denunciava havia anos a
pulverização aérea de agrotóxicos sobre os perímetros irrigados da região, onde
aviões sobrevoavam plantações de banana e melão vizinhas a casas, escolas e
açudes de abastecimento. A denúncia não tratava do veneno aplicado na lavoura,
mas do veneno que não fica na lavoura. Nove anos depois, o Ceará proibiu a
prática em todo o território e deu à lei o nome do homem assassinado (Lei
Estadual nº~16.820/2019).

A pulverização aérea é um caso de externalidade em que o mecanismo físico é
visível. Parte da calda liberada pela aeronave não atinge o alvo: é carregada
pelo vento e se deposita fora da área tratada, fenômeno que a literatura
técnica chama de deriva. Quem contrata a aplicação paga pelo serviço e pelo
produto; quem mora a favor do vento arca com o restante. A lei cearense
responde a esse conflito por proibição, e não se localizou, na literatura
consultada, avaliação causal dos efeitos dessa resposta.

\textbf{Qual é o efeito causal da proibição da pulverização aérea de
agrotóxicos sobre o peso ao nascer nos municípios cearenses onde havia
aplicação aérea?} É a pergunta deste projeto, e ela é de avaliação de política
pública. O objeto é a retirada de um \emph{método de aplicação}, e não de uma
substância, e essa distinção organiza o desenho inteiro: o grupo tratado
relevante é o dos municípios onde havia avião, e não o dos municípios onde
havia lavoura. A resposta buscada é um contraste, nos municípios com aplicação
aérea documentada entre 2007 e 2011: o efeito sobre as coortes cujo início de
gestação cairia na estação de pulverização, em relação às expostas a ela
apenas no fim da gestação. A estação de pulverização são os meses em que os voos se
concentram, provisoriamente a quadra chuvosa, de fevereiro a maio; o período
de referência, as coortes nascidas de 2007 a 2011.

\section{Antecedentes na literatura}

A evidência econômica próxima se divide entre o que uma substância faz à saúde
e o que um método de aplicação faz à saúde. Do lado da substância,
\citeonline{dias2023down} exploram a expansão da soja transgênica no Brasil,
combinada à direção do fluxo d'água dentro das bacias, e documentam piora dos
desfechos ao nascer nas populações a jusante, e
\citeonline{reynier2025glyphosate} associam a difusão do glifosato nos Estados
Unidos a menor peso ao nascer e gestação mais curta.

Do lado do método, os antecedentes diretos são dois.
\citeonline{camacho2017health} avaliam as consequências sanitárias da aspersão
aérea de herbicida sobre cultivos ilícitos na Colômbia.
\citeonline{calzada2023bananas} estudam a fumigação aérea de bananais no
Equador e encontram déficit de peso ao nascer entre filhos de mães que moravam
junto às plantações e foram expostas, na gestação, aos períodos de fumigação
intensa. Este último é o análogo mais próximo do caso cearense, pela cultura e
pelo método, com diferenças que importam: a doença combatida não é a mesma e,
pelos relatos disponíveis, a frequência das aplicações é menor
(Capítulo~\ref{cap:contexto}).

Entre os trabalhos localizados, nenhum avalia uma \textbf{proibição de
método}. Há simulações do banimento de moléculas específicas
\cite{bocker2020glyphosate}, mas simulação \emph{ex ante} não é avaliação
\emph{ex post}. Nos estudos de método, a pulverização aérea é introduzida ou
está em curso; no Ceará ela é retirada, e os sinais esperados se invertem.

\section{Justificativa}

A questão tem relevância prática imediata. A lei cearense foi contestada no
Supremo Tribunal Federal, que julgou constitucional a vedação em 2023
\cite{stf2023adi6137}, o que assentou a competência dos estados para legislar
sobre a matéria e torna a medida replicável; em 2024 o próprio Ceará
flexibilizou a regra para admitir drones, e a nova redação voltou a ser
questionada na Corte \cite{stf2025adi7794}. O debate tem se apoiado em
argumentos de plausibilidade, porque não se localizou estimativa causal do que
a proibição de fato produziu. Uma avaliação que diga se houve ganho em peso ao
nascer, ou que mostre com rigor por que ele não é mensurável com os registros
existentes, informa os estados que consideram adotar a medida e a decisão de
mantê-la, relaxá-la ou substituí-la. O alcance é limitado: o ensaio informa um
desfecho, por uma via, em nível municipal, e não mede os demais danos
atribuídos à pulverização.

\section{Contribuição e organização do projeto}

A contribuição pretendida tem três camadas. A primeira é de \textbf{objeto}: a
avaliação \emph{ex post} de uma proibição de método de aplicação. A segunda é
de \textbf{medida}: o projeto trata como problema central a distância entre a
cultura plantada e o método de aplicação efetivamente usado, e propõe medi-la
com registros administrativos; ela depende de os registros serem obtidos. A
terceira é de \textbf{identificação}: como a lei vale para todo o estado desde
a mesma data, o projeto usa um desenho de calendário gestacional
\cite{calzada2023bananas,brainerd2014seasonal}, com os municípios de banana
familiar como grupo de comparação e um período de referência fixado de
antemão, e o confronta com um desenho de intensidade de exposição
\cite{callaway2024continuous}. O teste só se realiza se houver voos
documentados e poder; em todos os cenários previstos, inclusive o de ele não
se realizar, o projeto diz o que entrega (Capítulo~\ref{cap:resultados}).

Este documento é o projeto do primeiro ensaio da dissertação; o segundo, sobre
a escolha do instrumento, é apenas mencionado e terá projeto próprio.
Descreve-se um plano de trabalho, não um resultado, embora duas rodadas
exploratórias já tenham cruzado a exposição com o desfecho do pós-lei
(Seção~\ref{sec:inferencia}). O Capítulo~\ref{cap:contexto} apresenta a lei e
a região; o Capítulo~\ref{cap:referencial}, a literatura; o
Capítulo~\ref{cap:objetivos}, os objetivos e as hipóteses; o
Capítulo~\ref{cap:metodologia}, a estratégia empírica; o
Capítulo~\ref{cap:dados}, os dados; e os
Capítulos~\ref{cap:resultados} e~\ref{cap:cronograma}, os resultados esperados
e o cronograma.
